% !TEX program = xelatex
% À compiler sur Overleaf avec le compilateur XeLaTeX
% (Menu -> Compiler -> XeLaTeX) pour avoir la police Verdana.

\documentclass[12pt,a4paper]{article}

\usepackage{fontspec}
% Verdana si disponible, sinon une police sans-serif proche
\IfFontExistsTF{Verdana}{\setmainfont{Verdana}}{\setmainfont{DejaVu Sans}}

\usepackage[french]{babel}
\usepackage[margin=1.4cm]{geometry}
\usepackage{amsmath}
\usepackage{setspace}
\usepackage{enumitem}
\usepackage{array}

% Interligne 1,5
\onehalfspacing

% Espacement réduit autour des égalités successives
\setlength{\abovedisplayskip}{1pt}
\setlength{\belowdisplayskip}{1pt}
\setlength{\abovedisplayshortskip}{0pt}
\setlength{\belowdisplayshortskip}{0pt}

% ----- Pas de numérotation des pages -----
\pagestyle{empty}

\setlength{\parindent}{0pt}

% ----- Commandes personnalisées -----
\newcounter{exercice}
\newcommand{\exercice}{%
  \refstepcounter{exercice}%
  \vspace{0.4em}
  \noindent\underline{\textbf{Exercice \arabic{exercice}}}\par
  \vspace{0em}
}

\setlist[enumerate,1]{label=\alph*),itemsep=0pt,topsep=1pt}
\setlist[itemize,1]{leftmargin=1.6em,itemsep=0pt,topsep=1pt}

\begin{document}

% ----- En-tête -----
\begin{center}
  \underline{\Large\textbf{Correction du contrôle séquence 13 version 1}}\\[0.3em]
 
\end{center}

% =====================================================================
\vspace{3cm}\exercice
\begin{enumerate}
  \item \textbf{Oui.} $36  = 6 \times 6$ 
  \newline 6 est donc un diviseur de 36.
  \item \textbf{Non.} $148 \div 7 \approx 21{,}14$
  \newline le quotient n'est pas un entier,
        donc 148 n'est pas un multiple de 7.
  \item \textbf{Non.}  $10 < 20$, donc 10 n'est pas un multiple de 20.
  \item \textbf{Non.} $38 \div 4 = 9{,}5$ 
  \newline le quotient n'est pas un entier,
        donc 4 n'est pas un diviseur de 38.
\end{enumerate}

% =====================================================================
\vspace{1cm}\exercice
\begin{enumerate}
  \item \textbf{2 et 3} (acceptés aussi : 5 ; 7 ; 11 ; \dots).
        Un nombre premier a exactement deux diviseurs : 1 et lui-même.
  \item \textbf{4 et 6} (acceptés aussi : 1 ; 8 ; 9 ; \dots).
      
  \item Dans 12 ; 16 ; 11 ; 25 ; 49 ; 56 ; 83 ; 93, les deux nombres premiers sont
        \textbf{11} et \textbf{83}.
       
\end{enumerate}

% =====================================================================
\vspace{1cm}\exercice
\textbf{a) } $6 = 2 \times 3$, 
donc les diviseurs de 6 sont $\boxed{1\ ;\ 6\ ;\ 2\ ;\ 3\ }$.

\textbf{b) } $20 = 2 \times 2 \times 5$,
donc les diviseurs de 20 sont
$\boxed{1\ ;\ 20\ ;\ 2\ ;\ 5\ ;\ 4\ ;\ 10\ }$.

% =====================================================================
\vspace{1cm}\exercice


\noindent\begin{minipage}[t]{0.48\textwidth}
\textbf{a)}
\begin{align*}
14 &= 2 \times 7
\end{align*}

\textbf{b)}
\begin{align*}
30 &= 2 \times 15\\
   &= 2 \times 3 \times 5
\end{align*}
\end{minipage}\hfill
\begin{minipage}[t]{0.48\textwidth}
\textbf{c)}
\begin{align*}
60 &= 2 \times 30\\
   &= 2 \times 2 \times 15\\
   &= 2 \times 2 \times 3 \times 5\\
\end{align*}

\textbf{d)}
\begin{align*}
105 &= 3 \times 35\\
    &= 3 \times 5 \times 7
\end{align*}
\end{minipage}

% =====================================================================
\exercice
\noindent\begin{minipage}[t]{0.31\textwidth}
\textbf{a)} 
\begin{align*}
\dfrac{3}{12} &= \dfrac{3 \times 1}{3 \times 4}\\
              &= \dfrac{1}{4}
\end{align*}
\end{minipage}\hfill
\begin{minipage}[t]{0.31\textwidth}
\textbf{b)} 
\begin{align*}
\dfrac{16}{24} &= \dfrac{8 \times 2}{8 \times 3}\\
               &= \dfrac{2}{3}
\end{align*}
\end{minipage}\hfill
\begin{minipage}[t]{0.31\textwidth}
\textbf{c)} 
\begin{align*}
\dfrac{60}{105} &= \dfrac{15 \times 4}{15 \times 7}\\
                &= \dfrac{4}{7}
\end{align*}
\end{minipage}

% =====================================================================
\vspace{1cm}\exercice
Le nombre de paquets doit diviser à la fois 105 et 60 : c'est un diviseur commun à
105 et 60. Pour en faire un maximum, on cherche le plus grand diviseur commun à 105 et 60.

Avec les décompositions en facteurs premiers obtenues à l'exercice 4 :
$105 = 3 \times 5 \times 7$ et $60 = 2 \times 2 \times 3 \times 5$.
Le plus grand diviseur commun à 105 et 60 est le produit des facteurs premiers communs :
3 \times 5
                      = 15

Hulk et Batman peuvent donc faire au maximum \textbf{15 paquets},
chacun contenant  7 cookies et 4 tartelettes.

% =====================================================================
\vspace{1cm}\exercice
On utilise les critères de divisibilité :
\newline \textbf{par 2} : le chiffre des unités est 0, 2, 4, 6 ou 8 ;
\newline \textbf{par 3} : la somme des chiffres est divisible par 3
(sommes des chiffres : $19$ pour 254\,350 ; $9$ pour 111\,111\,111 ; $22$ pour 811\,255) ;
\newline \textbf{par 5} : le chiffre des unités est 0 ou 5 ;
\newline \textbf{par 10} : le chiffre des unités est 0.

\renewcommand{\arraystretch}{1.4}
\setlength{\tabcolsep}{2pt}
\begin{center}
\small
\begin{tabular}{|c|c|c|c|c|}
\hline
 & \textbf{Divisible par 2} & \textbf{Divisible par 3} & \textbf{Divisible par 5} & \textbf{Divisible par 10} \\
\hline
254\,350 & Oui & Non & Oui & Oui \\
\hline
111\,111\,111 & Non & Oui & Non & Non \\
\hline
811\,255 & Non & Non & Oui & Non \\
\hline
\end{tabular}
\end{center}

\end{document}
